\begin{abstract}

\begin{center}

\large{Применение методов каузального вывода при анализе зависимостей между оптимизирующими проходами LLVM-компилятора} \\

\end{center}

\vspace{0.5cm}

Современные оптимизирующие компиляторы предоставляют десятки и сотни оптимизирующих проходов, эффект которых существенно зависит от порядка применения. Подбор оптимального порядка --- известная задача \emph{phase ordering problem} (далее в работе --- проблема поиска порядка применения оптимизаций) --- решается на практике эвристически (фиксированные уровни \texttt{-O2}, \texttt{-Oz}) либо методами машинного обучения, которые рассматривают компилятор как чёрный ящик и не объясняют, почему тот или иной порядок проходов оказывается выгодным. Существующие подходы анализируют корреляции между порядком проходов и итоговым качеством кода, что не позволяет отделить причинно-следственные зависимости от случайных совпадений.

В работе предлагается метод обнаружения каузальных зависимостей между оптимизирующими проходами LLVM-компилятора на основе рандомизированного эксперимента. Случайные перестановки набора из 89~проходов на выборке из 517~бенчмарков образуют рандомизированный эксперимент в смысле Рубина--Неймана: средний эффект порядка пары проходов (Average Treatment Effect) совпадает с её каузальным эффектом без необходимости устранения конфаундеров методами поиска структуры. На основе 5,17~миллиона прогонов построены каузальные графы зависимостей между проходами на трёх уровнях агрегации --- глобальный, по доменным классам программ и по кластерам, выделенным по структуре промежуточного представления.

Показано, что 92{,}7\% статистически значимых корреляционных рёбер между порядком пар проходов и метрикой улучшения кода не имеют каузальной природы; каузальный граф из 608~рёбер существенно беднее корреляционного из 5124~рёбер, построенного на тех же данных. Эффект порядка пар проходов гетерогенен по типам программ: до 29\% рёбер глобального графа меняют знак на отдельных кластерах. Дополнительно исследованы две формулировки воздействия --- относительный порядок пары проходов и их непосредственное соседство, --- выявляющие глобальные и локальные классы зависимостей. Собранные датасеты (517~скомпилированных бенчмарков, 5,17~миллиона прогонов с метриками) представляют самостоятельный научный вклад работы.

\end{abstract}

\newpage
